% ═══════════════════════════════════════
%   INTRODUCTION GÉNÉRALE
% ═══════════════════════════════════════
\chapter*{Introduction}
\addcontentsline{toc}{chapter}{Introduction}

Ce cahier de recherche documente mon apprentissage approfondi de dix piliers fondamentaux des mathématiques : la \textbf{logique et les raisonnements}, les \textbf{nombres complexes}, les \textbf{nombres réels}, les \textbf{suites numériques}, l'\textbf{arithmétique}, les \textbf{ensembles et applications}, les \textbf{limites et fonctions continues}, la \textbf{dérivée d'une fonction}, les \textbf{groupes}, et les \textbf{polynômes}. Il ne s'agit pas d'une simple recopie de cours, mais d'une reconstruction personnelle des concepts, enrichie de mes propres reformulations, intuitions et démonstrations détaillées.

\begin{intuitionbox}[Démarche]
Ma méthode d'apprentissage repose sur un principe directeur : pour chaque concept, je cherche la \emph{règle génératrice} sous-jacente plutôt que de mémoriser le résultat en surface. La question que je me pose systématiquement est :

\medskip
\centerline{\textit{<<~Quelle règle me permettrait de retrouver cela si je l'avais oublié~?~>>}}
\medskip

Ce cahier reflète cette approche. Chaque section reconstruit la théorie à partir de ses fondations, et mes commentaires personnels (dans des encadrés comme celui-ci) explicitent les liens structurels que j'ai identifiés.
\end{intuitionbox}

\textbf{Structure du cahier.} Chaque chapitre suit le même schéma :
\begin{enumerate}[leftmargin=2cm]
  \item \textbf{Définitions et constructions} --- les objets fondamentaux, posés rigoureusement.
  \item \textbf{Propriétés et théorèmes} --- les résultats clés avec démonstrations complètes.
  \item \textbf{Applications et exemples} --- mise en pratique concrète.
  \item \textbf{Synthèse personnelle} --- mes observations sur les connexions entre concepts.
\end{enumerate}
